\chapter{A Concrete Naming Certificate for $\EquicontinuityUp$}
\label{ch:concrete-instances-equicontinuity-namecert}
\origin{human}

Following Bishop and Bridges constructive analysis and Kelley-style compact-metric reasoning, the $\EquicontinuityUp$ packet records the finite family-modulus surface by which a compact source, a finite continuous-family ledger, and a shared tolerance-radius schedule are presented before any finite-net selector or Arzela--Ascoli consumer may read the family. It sits over the compact source of \autoref{ch:concrete-instances-compactmetric-namecert}, the continuous graph and pointwise-modulus rows of \autoref{ch:concrete-instances-continuousmap-namecert}, the tolerance-only modulus consumer of \autoref{ch:concrete-instances-uniformmodulus-namecert}, and the finite-net family route of \autoref{ch:concrete-instances-finite-net-equicontinuity-namecert}.

The packet exposes a shared family modulus. It does not assert compactness, select subsequences, provide a finite net, choose an Arzela--Ascoli limit, or import host-level continuity.

\begin{definition}[Equicontinuity carrier]
\label{def:equicontinuity-carrier}
An $\EquicontinuityUp$ carrier is a finite $\BHist$ packet
$$
E=(K,F,\epsilon,\rho,M,T,R,P,N)
$$
with the following displayed rows:
\begin{enumerate}
\item $K$ is the compact metric source row whose finite-cover and metric inspection data are read through $\CompactMetricUp$.
\item $F$ is the finite continuous-family row, each member of which is presented through $\ContinuousMapUp$ graph and pointwise-modulus data over the same source.
\item $\epsilon$ is the tolerance row consumed by every family member before a radius is exported.
\item $\rho$ is the shared radius row, read as a single family modulus rather than as one radius per function.
\item $M$ is the $\UniformEquicontinuityModulusUp$ or $\UniformModulusUp$ consumer row that records the shared tolerance-radius handoff.
\item $T$ records componentwise $\hsame$ transport for $K,F,\epsilon,\rho,M$.
\item $R$ records displayed $\Cont$ replay from compact-source observations through continuous-family rows into the shared-modulus handoff.
\item $P$ is the $\Pkg$ provenance over $\CompactMetricUp$, $\ContinuousMapUp$, $\UniformEquicontinuityModulusUp$, $\UniformModulusUp$, $\FiniteNetEquicontinuityUp$, $\BHist$, $\hsame$, $\Cont$, and $\NameCert$.
\item $N$ is the local $\NameCert_{\EquicontinuityUp}$ row.
\end{enumerate}
The carrier has no coordinate for a selected finite net, selected subsequence, compactness theorem, host continuity predicate, pointwise-only modulus family, or ambient completeness principle.
\end{definition}
\leandef{BEDC.Derived.EquicontinuityUp.EquicontinuityCarrier}

\begin{theorem}[Equicontinuity NameCert obligations]
\label{thm:equicontinuity-namecert-obligations}
For every accepted $\EquicontinuityUp$ carrier $E=(K,F,\epsilon,\rho,M,T,R,P,N)$, the local $\NameCert_{\EquicontinuityUp}$ surface consists exactly of carrier admission for the compact metric source, continuous-family ledger, tolerance row, shared radius row, modulus handoff, transport, replay, provenance, and local naming; shared-modulus exactness requiring the same $\rho$ to serve every member of $F$ at tolerance $\epsilon$; compact-source compatibility requiring any finite-net consumer to read $K$ before using the modulus row; and nonescape excluding compactness assertions, selected finite nets, selected subsequences, host continuity, and ambient completeness.
\end{theorem}
\leanchecked{BEDC.Derived.EquicontinuityUp.EquicontinuityCarrier\_namecert\_obligations}
\begin{proof}
Carrier admission is projection from \autoref{def:equicontinuity-carrier}. The tuple displays exactly $K,F,\epsilon,\rho,M,T,R,P,N$.

Shared-modulus exactness is the role of $\rho$ and $M$. The radius row is a single displayed row for the family, so a consumer cannot replace it by independently chosen radii for different members of $F$.

Compact-source compatibility is the route through $K$. A finite-net or compact-uniform consumer must read the compact metric source before it can consume the shared family modulus.

Finally, $T,R,P,N$ transport, replay, source, and name the same displayed rows. Since the carrier has no selected net, selected subsequence, compactness theorem, host-continuity predicate, or ambient-completeness coordinate, the listed obligations exhaust the local NameCert surface.
\end{proof}

\begin{theorem}[Equicontinuity compact-metric route]
\label{thm:equicontinuity-compact-metric-route}
Every compact-uniform or finite-net consumer of an accepted $\EquicontinuityUp$ carrier factors through
$$
K,F,\epsilon,\rho,M \longmapsto T,R,P,N .
$$
The route first reads the compact metric source $K$, then the continuous-family ledger $F$, then the shared tolerance-radius handoff $\epsilon,\rho,M$. It may be consumed by the finite-net equicontinuity surface of \autoref{ch:concrete-instances-finite-net-equicontinuity-namecert}, but it does not itself supply a finite net or compactness proof.
\end{theorem}
\begin{proof}
Project the accepted carrier. The only source coordinate is $K$, so any compact-uniform consumer starts from the displayed compact metric row.

Read the family row $F$ next. Its entries are continuous-map rows over the same source, so the family cannot change source data while the modulus is being consumed.

The shared handoff is $\epsilon,\rho,M$. It records a tolerance and one family radius before a downstream finite-net selector reads the packet. Transport and replay through $T,R$ preserve this order, while $P,N$ merely package and name it.
\end{proof}
\leanchecked{BEDC.Derived.EquicontinuityUp.EquicontinuityCompactMetricRoute}

\begin{theorem}[Equicontinuity nonescape]
\label{thm:equicontinuity-nonescape}
An accepted $\EquicontinuityUp$ carrier exports only the finite family-modulus rows $K,F,\epsilon,\rho,M,T,R,P,N$. It cannot export compactness, a selected finite net, a selected subsequence, Arzela--Ascoli convergence, host-level continuity, pointwise-only radius choice, or ambient $\RealUp$ completeness.
\end{theorem}
\begin{proof}
The carrier definition lists the entire public inventory. The analytic rows are the compact source, continuous family, tolerance, shared radius, and modulus handoff.

The structural rows transport and replay those same rows. They do not add a selector, a compactness principle, a subsequence object, a convergence row, or an external continuity predicate.

Thus every downstream route must obtain finite nets, compactness conclusions, and convergence statements from separate consumer chapters. The $\EquicontinuityUp$ packet supplies only the displayed family-modulus boundary.
\end{proof}
\leanchecked{BEDC.Derived.EquicontinuityUp.EquicontinuityTasteGate\_single\_carrier\_alignment}

\begin{theorem}[Equicontinuity family-window exposure]
\label{thm:equicontinuity-family-window-exposure}
For every accepted $\EquicontinuityUp$ carrier $E=(K,F,\epsilon,\rho,M,T,R,P,N)$, the family window exposed to compact-uniform consumers is the finite $\BHist$ route through the $\CompactMetricUp$ source row $K$, the $\ContinuousMapUp$ family ledger $F$, the tolerance row $\epsilon$, the shared radius row $\rho$, the $\UniformModulusUp$ handoff $M$, componentwise $\hsame$ transport $T$, displayed $\Cont$ replay $R$, $\Pkg$ provenance $P$, and local $\NameCert_{\EquicontinuityUp}$ row $N$. Each pointwise modulus read for a member of $F$ must occur inside that window before any finite-net radius is consumed.
\end{theorem}
\leanchecked{BEDC.Derived.EquicontinuityUp.EquicontinuityFamilyWindowExposure}
\begin{proof}
Project the carrier rows. The finite family window is exactly the displayed sequence $K,F,\epsilon,\rho,M,T,R,P,N$, so the compact source and family ledger are read before the shared-radius handoff.

The $\ContinuousMapUp$ pointwise-modulus entries live in $F$. They cannot be detached from the $\CompactMetricUp$ source $K$, because the compact-metric route forces every downstream read through the same source row.

Finally, $T,R,P,N$ only transport, replay, source, and name this window. A finite-net consumer therefore receives pointwise modulus data only through the displayed family window and not as an external radius choice.
\end{proof}

\begin{theorem}[Equicontinuity compact-net radius exactness]
\label{thm:equicontinuity-compact-net-radius-exactness}
Let $E=(K,F,\epsilon,\rho,M,T,R,P,N)$ be an accepted $\EquicontinuityUp$ carrier. Any $\CompactMetricUp$ finite-net radius consumed from $E$ is the same shared row $\rho$ that is handed to $\UniformModulusUp$ by $M$ after the finite $\ContinuousMapUp$ family window has been displayed. The $\BHist$, $\hsame$, $\Cont$, $\Pkg$, and $\NameCert$ rows preserve this radius exactly and add no memberwise radius ledger, selected net, or host compactness proof.
\end{theorem}
\leanchecked{BEDC.Derived.EquicontinuityUp.EquicontinuityCarrier\_compact\_net\_radius\_exactness}
\begin{proof}
Read $\rho$ between the tolerance row $\epsilon$ and the handoff row $M$. The carrier contains one shared radius row, not a list of member-indexed radii.

Use the compact-metric route and family-window exposure. The compact source $K$ and the finite continuous-family ledger $F$ are visible before the radius is consumed by a finite-net or uniform-modulus reader.

Transport and replay preserve the row rather than recomputing it. Thus the finite-net radius available to a compact consumer is exactly $\rho$ through $T,R,P,N$, with no selected net or host compactness coordinate added by the packet.
\end{proof}

\begin{theorem}[Equicontinuity uniform-modulus handoff]
\label{thm:equicontinuity-uniform-modulus-handoff}
Every $\UniformModulusUp$ handoff exported by an accepted $\EquicontinuityUp$ carrier factors through the finite route
$$
K,F,\epsilon,\rho,M,T,R,P,N .
$$
The handoff reads the $\CompactMetricUp$ finite-net radius boundary through $K,\rho$, reads each $\ContinuousMapUp$ pointwise modulus through the finite family row $F$, and then exposes $M$ as the shared $\UniformModulusUp$ consumer row. The $\BHist$, $\hsame$, $\Cont$, $\Pkg$, and $\NameCert$ support rows export no pointwise-only modulus family, selected finite net, or ambient $\RealUp$ completeness claim.
\end{theorem}
\leanchecked{BEDC.Derived.EquicontinuityUp.EquicontinuityUniformModulusHandoff}
\begin{proof}
Start from family-window exposure. It fixes the order in which $K$, $F$, $\epsilon$, and $\rho$ are displayed before the handoff row $M$ can be read.

Apply compact-net radius exactness. The radius consumed by the compact finite-net boundary is the same shared row $\rho$, and every pointwise modulus entry remains attached to the finite $\ContinuousMapUp$ ledger $F$.

The structural rows then repack the handoff. Since $T,R,P,N$ preserve and name only the displayed route, the exported $\UniformModulusUp$ row is a shared handoff and not a selected net, pointwise-only radius family, or completeness principle.
\end{proof}
\leanchecked{BEDC.Derived.EquicontinuityUp.EquicontinuityCarrier\_uniform\_modulus\_handoff}

\input{parts/concrete_instances/equicontinuity/obligation_closure_tail.tex}

\closureat{EquicontinuityUp}{scopedStr}

\begin{closurestatus}{\EquicontinuityUp}
  \constructivestory{An $\EquicontinuityUp$ packet is generated from a compact metric source row, a finite continuous-family ledger, shared tolerance and radius rows, a $\UniformEquicontinuityModulusUp$ or $\UniformModulusUp$ handoff row, componentwise hsame transport, displayed Cont replay, Pkg provenance, and local NameCert evidence.}
  \theoryclosure{\scopedClosure}
  \scopeclosed{\autoref{def:equicontinuity-carrier}, \autoref{thm:equicontinuity-namecert-obligations}, \autoref{thm:equicontinuity-compact-metric-route}, \autoref{thm:equicontinuity-nonescape}, \autoref{thm:equicontinuity-shared-radius-stability}, \autoref{thm:equicontinuity-shared-radius-obligation}, \autoref{thm:equicontinuity-finite-net-readiness}, and \autoref{thm:equicontinuity-modulus-ledger-exactness} seed the compact-source finite-family shared-modulus route over compact-source, continuous-family, tolerance, shared-radius, modulus-handoff, transport, replay, provenance, and local naming rows.}
  \formalstatus{\theoremCheckedV}
  \leantarget{BEDC.Derived.EquicontinuityUp.EquicontinuityRootObligationSourceFamily}
  \bridgestatus{none}
  \notclaimed{This chapter does not close compactness, selected finite nets, selected subsequences, Arzela--Ascoli convergence, host-level continuity, pointwise-only radius choice, ambient $\RealUp$ completeness, Lean verification, or bridge checking for $\EquicontinuityUp$.}
  \upgradepath{Public closure requires a public finite-net and Arzela--Ascoli-facing handoff that exports the scoped family-modulus route without adding selected nets, selected subsequences, or ambient $\RealUp$ completeness.}
\end{closurestatus}
